\section{Proof Patterns: Classical Form and Cred Form}
\label{sec:proof-patterns}

The proof layer should keep the old lesson from the material
conditional separate from the new conditioning mechanism.  A proof is
not a truth table entry.  It is a way of eliminating countermodels or
transporting bounds through verified constraints.  The central example
is the ordinary sequent
\[
  A\models B,
\]
read as ``$B$ follows from $A$'' or ``$B$ if $A$ is given.''  In
classical propositional logic this can be checked by a truth table, but
one need not compare the whole table symmetrically.  The only forbidden
row for $A\models B$ is
\[
  A=1,\qquad B=0.
\]
Equivalently, $A\models B$ says that $A\land\neg B$ is unsatisfiable,
or, classically, that $A\to B$ is valid.  The last formulation is useful
but dangerous as prose: it internalizes a metalinguistic consequence
relation as an object-language material conditional.

A safe informal reading is therefore:
\[
\boxed{\text{if $A$ is designated, then $B$ must be designated;
if $A$ is not designated, this sequent imposes no constraint.}}
\]
For the classical two-valued case, ``designated'' means true.  Thus
$A\models B$ checks only the rows in which $A$ is true.  Rows in which
$A$ is false are irrelevant for this particular sequent.  This is the
semantic reason why the material conditional can be used classically:
$A\to B$ is true exactly when the forbidden row $A=1,B=0$ is absent.
The danger is to read this internal formula as a genuine connection
between $A$ and $B$ rather than as a compact encoding of a countermodel
test.

The same pattern has three useful value-level forms:
\[
\begin{array}{@{}lll@{}}
\text{Boolean} & A,B\in\{0,1\}: & A\models B\iff
  \text{no valuation has }A=1,\ B=0.\\[1mm]
\text{Three-valued} & A,B\in\{0,\half,1\}: &
  \text{choose designated values, e.g. }\{\half,1\}\text{ for LP or }\{1\}\text{ for K}_3.\\[1mm]
\text{Continuous} & A,B\in[0,1]: & A\models_t B\iff
  t\le\val{A}_v\Rightarrow t\le\val{B}_v\text{ for every }v.
\end{array}
\]
In the three-valued case, the forbidden rows depend on the designation
policy.  LP forbids a designated antecedent with an undesignated
conclusion, so antecedent values $\half$ and $1$ both matter.  K$_3$
forbids only a certain antecedent with a non-certain conclusion, so only
antecedent value $1$ matters.  In the continuous case, the designation
policy is a threshold: one inspects all valuations where $A$ reaches the
chosen status $t$ and requires $B$ to reach the same status.

Cred keeps these sequent readings and rejects the material conditional
as a primitive.  For a threshold $t$, the sequent is
\[
  A\models_t B
  \quad\Longleftrightarrow\quad
  \forall v,\; t\le \val{A}_v \Rightarrow t\le \val{B}_v.
\]
The forbidden object is the graded countermodel
\[
  t\le \val{A}_v,
  \qquad
  \val{B}_v<t.
\]
Thus a Cred proof of $B$ from $A$ is still countermodel elimination,
but the countermodels are labelled by the demanded status: positive,
certain, or threshold $t$.  No formula $A\to B$ is introduced for this
purpose.

There are therefore three different statements that natural language
may blur:
\[
\begin{array}{ll}
\text{entailment:} & A\models_t B,\\[1mm]
\text{conditional credence:} & \Adm(c;j,e)\text{ with }e=\cred(A),\ j=\cred(A\land B),\\[1mm]
\text{internal implication:} & A\to B\text{ as a formula, not primitive in Cred.}
\end{array}
\]
Only the first is proof-theoretic consequence.  Only the second is
conditioning.  The third is deliberately absent from the core.

\begin{table}[h]
\centering
\begin{tabular}{@{}p{0.20\linewidth}p{0.34\linewidth}p{0.34\linewidth}@{}}
\toprule
Pattern & Classical reading & Cred reading \\
\midrule
Direct proof of $A\models B$ &
Assume $A$ and derive $B$.  Semantically, show that no valuation has
$A=1$ and $B=0$.  In a truth table, inspect only the rows where $A$ is
true and check that $B$ is true there. &
Assume the labelled premise $t\le\val{A}_v$ and derive the labelled
conclusion $t\le\val{B}_v$.  Equivalently, prove
$A\models_t B$ by showing that every valuation satisfying the premise
label also satisfies the conclusion label. \\
\addlinespace
Countermodel proof / disproof &
To refute $A\models B$, exhibit one row/model with $A=1$ and $B=0$. &
To refute $A\models_t B$, exhibit one valuation with
$t\le\val{A}_v$ and $\val{B}_v<t$.  The countermodel may be graded; it
need not be Boolean. \\
\addlinespace
Indirect proof / reductio &
Assume $A$ and $\neg B$ and derive contradiction.  In model terms,
show that $A\land\neg B$ has no model. &
Assume the graded countermodel conditions
$t\le\val{A}_v$ and $\val{B}_v<t$ and show they are impossible.  Reductio
is legitimate because the countermodel set is empty, not because
contradiction itself licenses arbitrary conclusions. \\
\addlinespace
Contrapositive &
Use $A\models B$ iff $\neg B\models\neg A$, relying on classical
negation and the classical consequence relation. &
Use only when the corresponding semantic bridge is available.  Cred has
involutive negation as a value operation, but it does not add an
internal implication merely to make contraposition syntactic. \\
\addlinespace
Tautology check &
Check that $A\to B$ is true on every row.  This is classically
equivalent to $A\models B$. &
Avoid this as the primitive proof form.  The primitive proof form is the
external sequent $A\models_t B$ or a labelled derivation.  The absence
of an internal arrow prevents the material-conditional reading from
being mistaken for dependence. \\
\addlinespace
Explosion &
From inconsistent premises $A,\neg A$, every $B$ follows, because no
classical valuation satisfies both premises. &
Explosion is blocked at the relevant thresholds by surviving
countermodels.  A contradictory pair can be designated without forcing
an arbitrary conclusion. \\
\addlinespace
Case split &
Prove the conclusion under each case and combine the cases by excluded
middle or a partition of possibilities. &
Split the value/fiber space.  Positive evidence gives a singleton
fiber; zero evidence gives the whole interval; incoherent pairs give
empty fibers.  Subsequent rules map or intersect the remaining
admissible set. \\
\addlinespace
Modus ponens &
From $A$ and $A\to B$, infer $B$, where $\to$ is often internalized as
a connective. &
There is no internal conditional formula.  The chain rule is used as a
side judgment: from evidence above threshold $t$ and an admissible
conditional bounded below by $\ell$, infer the joint above
$\ell\conjc t$. \\
\addlinespace
Bayesian update &
Given $P(A\cap B)$ and $P(B)>0$, divide to get $P(A\mid B)$. &
Use the admissible-set trichotomy.  Positive evidence gives the unique
conditional $j/e$; zero evidence carries the non-singleton fiber rather
than dividing by zero. \\
\addlinespace
Independence / irrelevance &
If $P(A\cap B)=P(A)P(B)$ and $P(B)>0$, then $P(A\mid B)=P(A)$: the
evidence does not update the conclusion. &
The joint is supplied dependence data.  A product joint gives trivial
conditioning; a non-product joint gives nontrivial update.  Truth values
alone do not determine either case. \\
\bottomrule
\end{tabular}
\caption{Classical proof patterns and their Cred counterparts.  The
Cred column replaces internal implication by labelled consequence,
countermodel elimination, and external conditioning judgments.}
\label{tab:proof-patterns}
\end{table}

\subsection{Where Cred Is More Elegant}

The valuable differences are not obtained by rejecting classical proof.
They come from making hidden bookkeeping explicit and local.  The most
important differences are the following.

\paragraph{1. Entailment is stated directly, not via a material arrow.}
Classically, one often proves $A\models B$ by proving the tautology
$A\to B$.  This is correct but invites a misleading reading of $A\to B$
as a connection.  Cred writes the sequent itself:
\[
  A\models_t B.
\]
The proof obligation is exactly the absence of a graded countermodel
$t\le\val{A}_v$ and $\val{B}_v<t$.  This is more transparent because the
object being ruled out is visible in the statement.

\emph{Example.}  To show $A\models B$ classically, one checks that the
row $A=1,B=0$ cannot occur.  To show $A\models_t B$ in Cred, one checks
that no valuation has $A$ at least at status $t$ while $B$ falls below
$t$.  The same proof idea works for Boolean truth, LP positivity, K$_3$
certainty, and arbitrary thresholds.

\paragraph{2. Proof status is labelled instead of hidden in prose.}
Classical mathematics usually says simply ``assume $A$'' or ``prove
$B$.''  Cred records whether the proof transports positivity, certainty,
or a threshold.  This prevents accidental changes of strength.

\emph{Example.}  Positivity consequence asks for $\val{A}_v>0\Rightarrow
\val{B}_v>0$.  Certainty consequence asks for $\val{A}_v=1\Rightarrow
\val{B}_v=1$.  These are different proof obligations.  Cred keeps them
separate, while the Kleene collapse explains them as LP and K$_3$
consequence.

\paragraph{3. Reductio is countermodel elimination, not explosion.}
In classical logic, reductio and explosion are close because an
inconsistent premise set has no models.  Cred separates the legitimate
part from the dangerous part.  Reductio is valid when the countermodel
set for the target is empty.  A contradiction among premises is not by
itself a license to infer everything.

\emph{Example.}  To prove $A\models_t B$ indirectly, assume a graded
countermodel: $t\le\val{A}_v$ and $\val{B}_v<t$.  Derive impossibility.
This proves the sequent.  But if both $A$ and $\negc A$ are merely
designated at a paraconsistent threshold, Cred may still have valuations
where $B$ fails.  Thus explosion is blocked because a countermodel
survives.

\paragraph{4. Conditioning is a side judgment, not an implication formula.}
Classical-looking proofs often blur ``$B$ follows from $A$'' with
``the conditional $B\mid A$ has value $1$.''  Cred keeps these apart.
Entailment is a statement over all valuations.  Conditioning is the
external admissibility equation
\[
  \Adm(c;j,e)\quad\Longleftrightarrow\quad c e=j.
\]

\emph{Example.}  If $e>0$, the equation gives the unique value
$c=j/e$.  If $e=j=0$, the equation gives every $c\in[0,1]$.  Thus zero
evidence does not prove every conclusion; it fails to determine a
conditional value.

\paragraph{5. Case splits become fiber propagation.}
Classical case splits usually collapse each branch to true or false.
Cred can instead carry the remaining admissible set.

\emph{Example.}  For positive evidence, the fiber is the singleton
$\{j/e\}$.  For zero evidence and zero joint, the fiber is $[0,1]$.
For incoherent data $j>e$, the fiber is empty.  Later reasoning maps,
intersects, or preserves these sets.  This is more informative than
choosing an arbitrary default value at the singular case.

\paragraph{6. Dependence is supplied explicitly.}
Classical truth tables only see the values of $A$ and $B$ at a row.
They do not encode whether $A$ and $B$ are probabilistically or
semantically connected.  Cred makes the probability analogy explicit:
the marginals do not determine the joint.

\emph{Example.}  If the supplied joint is the product joint
$j=\cred(A)\cred(B)$ and $\cred(B)>0$, then
\[
  \cred(A\mid B)=\cred(A),
\]
so the evidence does not update the conclusion.  If the joint is not
the product joint, the same marginal truth values may yield a nontrivial
update.  Thus coincident truth values never manufacture a connection.

Together, these differences make the proof language cleaner.  A
classical proof often works by encoding consequence as a material
conditional and relying on the reader to remember the semantic
translation.  A Cred proof states the semantic translation directly:
which valuations are being ruled out, which status is being transported,
which dependence data are supplied, and whether conditioning is uniquely
determined or only a fiber.

The conceptual improvement is not that classical proof techniques are
invalid.  Direct proof, case split, reductio, and threshold consequence
remain ordinary semantic methods.  The improvement is that Cred states
what each method is actually using.  Direct proof transports a labelled
bound from premises to conclusion.  Reductio eliminates the
countermodel set of the target.  Bayesian update inverts a chain-rule
constraint only when that inverse is single-valued.  At zero evidence,
the proof object is not an arbitrary conclusion but a propagated set of
admissible values.

This gives a clean replacement for the misleading material-conditional
reading of proof.  Classical mathematics may continue to use material
implication internally when appropriate.  Cred keeps a stricter
interface: object-language formulas have values; consequence is an
external relation over valuations; conditioning is an external
admissibility judgment; dependence is supplied by the joint.  This is
why the system can preserve the useful classical proof patterns without
turning false evidence into a universal premise.
